%! AP Statistics — Unit 1, Lesson 1.5 — Warm-Up ANSWER KEY
\documentclass[10pt]{article}
\usepackage{apstats-article}
\usepackage{apstats-key}

%=============================================================================
\begin{document}

\pageheader{Unit 1, Lesson 1.5}{Warm-Up --- Answer Key}
\begin{tabularx}{\linewidth}{X X X}
Name:\;\ans{Key} & Date:\; & Period:\;
\end{tabularx}

\vspace{0.16in}

\noindent The dotplot below shows the number of books read over the summer by
each of 20 students.

\vspace{0.14in}

\begin{center}
\begin{tikzpicture}[scale=1.0]
  \draw[thick] (0,0) -- (5.2,0);
  \foreach \x/\lbl in {0.8/1, 1.8/2, 2.8/3, 3.8/4, 4.8/5} {
    \draw (\x,-0.12) -- (\x,0.12);
    \node[below, font=\small] at (\x,-0.12) {\lbl};
  }
  \node[below, font=\small\itshape] at (2.8,-0.52) {Number of books read};
  \foreach \k in {1,...,4} \filldraw[navy] (0.8, \k*0.38) circle (0.14);
  \foreach \k in {1,...,7} \filldraw[navy] (1.8, \k*0.38) circle (0.14);
  \foreach \k in {1,...,6} \filldraw[navy] (2.8, \k*0.38) circle (0.14);
  \foreach \k in {1,...,2} \filldraw[navy] (3.8, \k*0.38) circle (0.14);
  \filldraw[navy] (4.8, 0.38) circle (0.14);
\end{tikzpicture}
\end{center}

\vspace{0.18in}

\begin{enumerate}[leftmargin=1.5em, itemsep=10pt]

  \item \ansline{The variable is ``number of books read over the summer.'' It is
        \textbf{quantitative} because the values are counts (numbers on a meaningful
        scale) --- you can add, subtract, and compute averages.}

  \item \ansline{The value \textbf{2} occurs most often. \textbf{7} students chose
        that value (7 dots stacked above the 2).}

  \begin{teachernote}
  Common error: students may say 3 (there are 6 dots above 3, which is close).
  Have students count each stack carefully.
  \end{teachernote}

  \item \ansline{Range $= 5 - 1 = \mathbf{4}$ books.}

  \item \ansline{In a dotplot, the number line has a meaningful numerical order and
        equal spacing --- the distance between 1 and 2 equals the distance between
        4 and 5. In a bar chart for categorical data, the order of bars is arbitrary
        and there is no numerical scale. Also, histogram bars (for quantitative data)
        \textit{touch} because adjacent bins share a boundary, whereas bar chart bars
        do not touch.}

  \begin{teachernote}
  Accept any one of: (1) the axis has a meaningful numerical scale; (2) order matters
  in a dotplot but not in a bar chart; (3) dotplot bars (stacks) would touch because
  adjacent values are consecutive on the number line. Do not require the student to
  mention histograms specifically --- that comes later in today's lesson.
  \end{teachernote}

\end{enumerate}

\end{document}
